% Part: incompleteness
% Chapter: arithmetization-syntax
% Section: proofs-in-ax

\documentclass[../../../include/open-logic-section]{subfiles}

\begin{document}

\olfileid{inc}{art}{pax}
\olsection{\usetoken{P}{derivation} Aksiomatik}

\begin{explain}
Supaya bisa ngaritmetisasi !!{derivation}s aksiomatik, kita kudu
makili !!{derivation}s minangka wilangan. Amarga !!{derivation}s
iki mung urutan !!{formula}s, cara kang cetha yaiku ngodeake
saben !!{derivation} minangka kode saka urutan kode !!{formula}s
ing njerone.
\end{explain}

\begin{defn}
Yen $\delta$ iku !!{derivation} aksiomatik kang dumadi saka
!!{formula}s $!A_1$, \dots,~$!A_n$, $\Gn{\delta}$ yaiku
\[
\tuple{\Gn{!A_1}, \dots, \Gn{!A_n}}.
\]
\end{defn}

\begin{ex}
  Gatekna !!{derivation} kang prasaja iki:
  \begin{derivation}
  1. & $!B \lif (!B \lor !A)$ \\
  2. & $(!B \lif (!B \lor !A)) \lif (!A  \lif (!B \lif (!B \lor !A)))$\\
  3. & $!A  \lif (!B \lif (!B \lor !A))$
  \end{derivation}
  Wilangan G\"odel saka !!{derivation} iki yaiku
  \begin{align*}
  \openTuple\, 
    & \Gn{!B \lif (!B \lor !A)}, \\
    &\Gn{(!B \lif (!B \lor !A)) \lif (!A  \lif (!B \lif (!B \lor !A)))},\\
    & \Gn{!A  \lif (!B \lif (!B \lor !A))} \,\closeTuple.
  \end{align*}
\end{ex}

\begin{explain}
Sawise nemtokake representasi !!{derivation}s, kita uga kudu
nuduhake yen !!{derivation}s iki bisa diolah kanthi rekursif
primitif, lan sifat sarta relasi pokoke bisa dinyatakake kanthi
cara kang padha. Sawetara operasi prasaja: upamane, yen diwenehi
wilangan G\"odel~$d$ saka !!a{derivation},
$(d)_{\len{d}-1}$ menehi wilangan G\"odel saka !!{formula}
pungkasan. Sawetara sifat liyane luwih angel. Kita bakal menehi
sketsa carane. Tujuane yaiku nuduhake yen relasi
``$\delta$ iku !!a{derivation} saka~$!A$ adhedhasar~$\Gamma$''
iku rekursif primitif ing wilangan G\"odel $\delta$ lan~$!A$.
\end{explain}

\begin{prop}
\ollabel{prop:followsby}
Relasi-relasi iki rekursif primitif:
\begin{enumerate}
\item $!A$ iku aksioma.
\item Baris kapangka-$i$ ing $\delta$ dibenerake dening modus ponens.
\item Baris kapangka-$i$ ing $\delta$ dibenerake dening \QR.
\item $\delta$ iku !!{derivation} kang bener.
\end{enumerate}
\end{prop}

\begin{proof}
Kita kudu nuduhake yen relasi kang cocog ing antarane wilangan
G\"odel saka !!{formula}s lan wilangan G\"odel saka
!!{derivation}s iku rekursif primitif.
\begin{enumerate}
\item Kita duwe dhaptar skema aksioma, lan $!A$ iku aksioma yen
  wujude cocog karo salah sijine skema mau. Amarga dhaptar skema
  winates, cukup nuduhake yen kanggo saben skema aksioma,
  kita bisa mriksa kanthi rekursif primitif apa $!A$ nduweni
  wujud kasebut. Upamane, gatekna skema aksioma
  \[
  !B \lif (!C \lif !B).
  \]
  $!A$ iku instans saka skema aksioma iki yen ana !!{formula}s
  $!B$ lan $!C$ saengga kita entuk $!A$ kanthi nyambungake `$($'
  karo $!B$, karo `$\lif$', karo `$($', karo $!C$, karo `$\lif$',
  karo $!B$, lan karo~`$))$'. Sifat kang cocog saka wilangan
  G\"odel $n$ saka~$!A$ bisa dipriksa, amarga panyambungan urutan
  iku rekursif primitif lan wilangan G\"odel saka $!B$ lan $!C$
  kudu luwih cilik tinimbang wilangan G\"odel saka~$!A$.
  Yen relasi iki bener, $!B$ lan $!C$ iku sub-!!{formula}s
  saka~$!A$. Mula kita bisa netepake:
  \begin{multline*}
  \fn{IsAx}_{!B \lif (!C \lif !B)}(n) \defiff
    \bexists{b< n}{\bexists{c < n}{
      \fn{Sent}(b) \land \fn{Sent}(c) \land {}\\
      n = \Gn{(} \concat b \concat \Gn{\lif} \concat \Gn{(}
      \concat c \concat \Gn{\lif} \concat b \concat \Gn{))}}}.
  \end{multline*}
  Yen saben skema aksioma nduweni definisi kaya mangkono,
  disjungsi definisi-definisi mau nemtokake sifat $\fn{IsAx}(n)$,
  yaiku ``$n$~iku wilangan G\"odel saka aksioma.''
\item Baris kapangka-$i$ ing $\delta$ dibenerake dening modus
  ponens yen lan mung yen ana baris $j$ lan $k < i$ kang
  !!{sentence} ing baris~$j$ iku sawijining formula $!A$,
  ukara ing baris~$k$ iku $!A \lif !B$, lan ukara ing
  baris~$i$ iku~$!B$.
  \begin{multline*}
    \fn{MP}(d, i) \defiff \bexists{j < i}{\bexists{k < i}{
      (d)_k = \Gn{(} \concat (d)_j
      \concat \Gn{\lif} \concat (d)_i \concat \Gn{)}}}
  \end{multline*}
Cathetan panarikan OLPL-182 lan OLPL-183: klaim cakupan lawas ditarik. Nutup argumen TeX ora nutup cakupan logis; rumus sumber wis ngiket saksi kang digunakake. Tata brace nested kang padha maknane dijaga.
  Amarga kuantifikasi winates, panyambungan, lan $=$ iku
  rekursif primitif, iki netepake relasi rekursif primitif.
\item Sawijining baris ing $\delta$ dibenerake dening \QR{} yen
  wujude $!B \lif \lforall[x][!A(x)]$, ana baris sadurunge
  kang awujud $!B \lif !A(c)$ kanggo sawijining !!{constant}~$c$,
  lan $c$ ora ana ing~$!B$. Iki lumaku yen lan mung yen
  \begin{enumerate}
  \item ana !!a{sentence}~$!B$ lan
  \item !!a{formula}~$!A(x)$ kang variabel bebas, yen ana, mung~$x$,
    saengga
  \item baris $i$ ngemot $!B \lif \lforall[x][!A(x)]$,
  \item sawijining baris $j < i$ ngemot $!B \lif \Subst{!A}{c}{x}$
    kanggo sawijining konstanta~$c$,
  \item kang ora ana ing~$!B$.
  \end{enumerate}
  Kabeh iki bisa dipriksa kanthi rekursif primitif, amarga wilangan
  G\"odel saka $!B$, $!A(x)$, lan $x$ luwih cilik tinimbang
  wilangan G\"odel saka formula ing baris~$i$, lan wilangan
  kode saka $c$ luwih cilik tinimbang wilangan G\"odel
  saka formula ing baris~$j$:
  \begin{multline*}
    \fn{QR}_1(d, i) \defiff
    \bexists{j<i}{\bexists{b < (d)_i}{\bexists{x < (d)_i}{
      \bexists{a < (d)_i}{\bexists{c < (d)_j}{
        \fn{Frm}(a) \land \fn{Var}(x) \land \fn{Const}(c) \land {}\\
        (d)_i = \Gn{(} \concat b \concat \Gn{\lif}
        \concat \Gn{\lforall} \concat x \concat a
        \concat \Gn{)} \land {}\\
        (d)_j = \Gn{(} \concat b \concat \Gn{\lif}
        \concat \fn{Subst}(a,c,x) \concat \Gn{)} \land {}\\
        \fn{Sent}(b) \land \fn{Sent}(\fn{Subst}(a,c,x)) \land {}
        \bforall{k < \len{b}}{(b)_k \neq (c)_0}}}}}}
  \end{multline*}
Cathetan koreksi sumber OLPL-184: saksi baris sadurunge j kudu dikuantifikasi, lan rumus Jawa wis ngenalake j kurang i. Bagean klaim lawas ngenani nutup argumen TeX ditarik.
Cathetan koreksi sumber OLPL-185: wates saka baris j ing tes QR ditrapake marang kode konstanta c, dene c lan x iku kode suku konstanta lan variabel miturut urutane; target wis nyocogake panganggone.
  Ing kene kita nganggep $c$ lan $x$ minangka wilangan G\"odel
  saka konstanta lan variabel kang dianggep suku (dudu kode
  simbule). Kita mriksa yen $x$ iku siji-sijine variabel bebas
  ing~$!A(x)$ kanthi mriksa apa $\Subst{!A(x)}{c}{x}$ iku
  !!a{sentence}, lan mesthekake yen $c$ ora ana ing $!B$
  kanthi mbutuhake saben simbol ing~$!B$ beda karo~$(c)_0$.

  Versi liyane saka \QR{} kita pasrahake dadi gladhen.
\item $d$ iku wilangan G\"odel saka !!{derivation} kang bener
  yen lan mung yen saben barise iku aksioma, utawa dibenerake
  dening modus ponens utawa~\QR. Mula:
  \begin{multline*}
  \fn{Deriv}(d) \defiff \fn{Seq}(d) \land {}\\
  \bforall{i < \len{d}}{(\fn{Sent}((d)_i) \land
    [\fn{IsAx}((d)_i) \lor \fn{MP}(d,i) \lor \fn{QR}(d, i)])}
  \end{multline*}
Cathetan koreksi sumber SOL6-F081: Seq nguji kode urutan sah, yaiku nol utawa siji kanggo kode kosong, utawa produk pangkat positif prima saka kapisan tanpa bolong. Iki rekursif primitif. Saben baris dicek ukara supaya modus ponens ora nampa gabungan kang sintakse salah.
\end{enumerate}
\end{proof}

\begin{prob}
Tetepna relasi-relasi iki kaya ing
\olref[inc][art][pax]{prop:followsby}:
\begin{enumerate}
\item $\fn{IsAx}_{!A \lif (!B \lif (!A \land !B))}(n)$,
\item $\fn{IsAx}_{\lforall[x][!A(x)] \lif !A(t)}(n)$,
\item $\fn{QR}_{2}(d, i)$ (kanggo versi liyane saka \QR).
\end{enumerate}
\end{prob}

\begin{prop}
Upamane $\Gamma$ iku himpunan rekursif primitif saka
!!{sentence}s. Relasi $\Prf[\Gamma](x, y)$ kang nyatakake
``$x$ iku kode saka !!a{derivation}~$\delta$ kanggo kondisional kang disusun nengen kanthi konsekuen $!A$
lan urutan anteseden winates saka~$\Gamma$, kang diturunake tanpa asumsi, lan $y$ iku wilangan G\"odel saka~$!A$''
iku rekursif primitif.
\end{prop}

\begin{proof}
Upamane ``$y \in \Gamma$'' diwenehake dening predikat
rekursif primitif~$R_\Gamma(y)$. Kita kudu nuduhake yen relasi
$\Prf[\Gamma](x, y)$ iku rekursif primitif; relasi
$\Prf[\Gamma](x, y)$ bener yen lan mung yen $y$ iku
wilangan G\"odel saka !!a{sentence}~$!A$ lan $x$~iku
kode saka !!a{derivation} tanpa asumsi kanggo kondisional kang disusun nengen kanthi konsekuen~$!A$ lan urutan anteseden winates saka $\Gamma$.

Miturut proposisi sadurunge, sifat $\fn{Deriv}(x)$, kang bener
yen lan mung yen $x$ iku kode saka !!{derivation}~$\delta$
kang bener, iku rekursif primitif. Nanging definisi mau durung
nggatekake himpunan $\Gamma$ minangka cara tambahan kanggo
mbenerake baris ing !!{derivation}. Tes rekursif primitif kita
apa sawijining baris dibenerake dening \QR{} uga durung
nggatekake syarat yen konstanta~$c$ ora kena ana ing~$\Gamma$.
Definisi bisa didandani supaya langsung nggatekake $\Gamma$,
nanging luwih gampang nggunakake $\fn{Deriv}$ lan teorema
deduksi. $\Gamma \Proves !A$ bener yen lan mung yen ana dhaptar
winates !!{sentence}s $!B_1$, \dots, $!B_n \in \Gamma$
saengga $\{!B_1, \dots, !B_n\} \Proves !A$. Miturut teorema
deduksi, iki lumaku yen $\Proves (!B_1 \lif (!B_2 \lif
\cdots (!B_n \lif !A)\cdots))$. Apa !!a{sentence} kanthi
wilangan G\"odel~$z$ awujud kaya iki bisa dipriksa kanthi
rekursif primitif. Mula, tinimbang nganggep $x$ minangka
wilangan G\"odel saka !!a{derivation} kanggo !!{sentence}
kanthi wilangan G\"odel~$y$ \emph{adhedhasar $\Gamma$},
kita nganggep $x$ minangka wilangan G\"odel saka
!!a{derivation} kanggo kondisional kang disusun nengen kaya ing ndhuwur
adhedhasar~$\emptyset$.

Pisanan, yen kita duwe urutan !!{sentence}s, kanthi rekursif
primitif kita bisa mbentuk kondisional kanthi kabeh ukara mau
minangka anteseden lan !!{sentence} kang diwenehake minangka
konsekuen:
\begin{align*}
  \fn{hCond}(s, y, 0) & = y \\
  \fn{hCond}(s, y, n+1) & = \Gn{(} \concat (s)_{\len{s} \tsub (n+1)} \concat \Gn{\lif}
  \concat \fn{hCond}(s, y, n) \concat \Gn{)}\\
  \fn{Cond}(s, y) & = \fn{hCond}(s, y, \len{s})\\
  \intertext{Mula kita bisa netepake $\Prf[\Gamma](x, y)$ dening}
  \Prf[\Gamma](x, y) & \defiff \fn{Sent}(y) \land \len{x}>0 \land {}\\
  &\bexists{s < \fn{sequenceBound}(x,x)}{(\fn{Seq}(s) \land {}} \\
  &\qquad (x)_{\len{x}-1} = \fn{Cond}(s,y) \land {} \\
  &\qquad\bforall{i<\len{s}}{(s)_i \in \Gamma} \land {} \\
  &\qquad\fn{Deriv}(x)).
\end{align*}
Cathetan koreksi sumber OLPL-186: langkah rekursi hCond ing sumber salah nyebut Cond kanthi telung argumen; rumus Jawa wis nganggo hCond kang cocog karo definisi rekursi.
Wates kanggo~$s$ diduweni kanthi nggatekake yen saben $(s)_i$
iku wilangan G\"odel saka sub-!!{formula} ing baris pungkasan
saka !!{derivation}, yaiku luwih cilik tinimbang
$(x)_{\len{x}-1}$. Cacah anteseden $!B \in \Gamma$, yaiku
dawane~$s$, luwih cilik tinimbang dawa baris pungkasan
ing~$x$.
Cathetan koreksi sumber OLPL-187 lan SOL6-F082: pratelan relasi lan panjelasan input saiki diowahi dadi sertifikat kondisional tanpa asumsi. Teorema deduksi menehi korespondensi anane bukti, dudu identitas kode bukti asli. SOL6-F080 ngolah urutan anteseden saka pungkasan supaya kondisional njaga urutan dhaptar. SOL6-F081 nguji kode urutan anteseden sah, ukara target lan derivasi ora kosong.
\end{proof}

\end{document}
